% !TEX program = xelatex
% csp-s-2026-answer.tex
% 2026 CCF 非专业级别软件能力认证第一轮（CSP-S1）提高级 C++语言试题（模拟卷）· 参考答案与解析
% 标准解析文档：每题给出答案、解析、考察知识点（NOI 大纲 2025）、难度系数、类别
% 编译方式：xelatex csp-s-2026-answer.tex（建议编译两次）
\documentclass[a4paper]{ctexart}

\usepackage{fontspec}
\usepackage{geometry}
\geometry{top=2.7cm, left=1.9cm, right=1.9cm, bottom=2.9cm}
\setmainfont{Consolas}
\setCJKmainfont{SimSun}
\newCJKfontfamily{\dengxian}{DengXian}

\usepackage{amsmath}
\usepackage{booktabs}
\usepackage{longtable}
\usepackage{array}
\usepackage{enumitem}

\setlength{\parindent}{0pt}
\setlength{\parskip}{3pt}
\linespread{1.25}

% 每题解答：题号+答案行，然后解析
\newcommand{\ans}[2]{\noindent\textbf{#1} \quad 答案：#2\par}
\newcommand{\expn}[1]{\hspace*{1.5em} 解析：#1\par}

\newcommand{\bigtitle}[1]{\begin{center}\dengxian\large #1\end{center}}

\begin{document}

\begin{center}
{\dengxian
    {\xeCJKsetup{CJKecglue={\hskip 0.419em}}\fontsize{15.95}{1}\selectfont 2026 CCF非专业级别软件能力认证第一轮}\\[1em]
    {\xeCJKsetup{CJKecglue={\hskip 0.419em}}\fontsize{14.05}{1}\selectfont （CSP-S1）提高级 C++语言试题 \textbf{参考答案与解析}}
}
\end{center}

\noindent 说明：类别分为基础知识与编程环境、C++ 程序设计、数据结构、算法、图论、数学与其他；难度系数与组别取自《全国青少年信息学奥林匹克系列竞赛大纲（2025 年修订版）》。

\section*{一、单项选择题（1--15）}

\ans{1.}{B}
\expn{\texttt{cat}（concatenate）用于连接并显示文件内容；\texttt{cp} 复制文件，\texttt{ping} 测试网络连通性，\texttt{touch} 创建空文件或更新时间戳。\quad 知识点：Linux 终端常用命令；类别：基础知识与编程环境；难度：5；组别：提高级。}

\ans{2.}{C}
\expn{哈夫曼树没有度为 1 的结点，\texttt{m} 个叶子对应 \texttt{2m-1} 个结点；二叉链表共 \texttt{2(2m-1)} 个指针域，其中非空指针数等于边数 \texttt{2m-2}，故空指针域为 \texttt{2(2m-1)-(2m-2) = 2m}。\quad 知识点：哈夫曼树性质；类别：数据结构；难度：4；组别：入门级。}

\ans{3.}{B}
\expn{完全二叉树中每个叶子到根的路径长度为 \texttt{O(log n)}，叶子约有 \texttt{n/2} 个，总访问量为 \texttt{O(n log n)}。\quad 知识点：完全二叉树与时间复杂度；类别：算法；难度：6；组别：提高级。}

\ans{4.}{D}
\expn{全程 \texttt{x >= y} 等价于任意前缀中向右步数不少于向上步数，方案数为卡特兰数 \texttt{C5 = 42}。\quad 知识点：卡特兰数；类别：数学与其他；难度：7；组别：提高级。}

\ans{5.}{A}
\expn{双栈模拟队列：出队时若栈 B 为空，先把栈 A 全部倒入栈 B（元素顺序反转），再弹出栈 B 栈顶；若栈 B 非空则直接弹出栈 B 栈顶，即为最早入队的元素。\quad 知识点：栈与队列的模拟；类别：数据结构；难度：3；组别：入门级。}

\ans{6.}{B}
\expn{设 \texttt{f(n)} 为登上 n 级台阶的走法数，最后一步跨 1 级或 2 级，故 \texttt{f(n) = f(n-1) + f(n-2)}。\texttt{f(1)=1}，\texttt{f(2)=2}，依次得 \texttt{f(10)=89}。\quad 知识点：递推；类别：算法；难度：3；组别：入门级。}

\ans{7.}{C}
\expn{\texttt{std::set} 中的元素按从小到大自动有序且不重复，插入、删除、查找的复杂度均为 \texttt{O(log n)}；\texttt{std::map} 按键有序存储且键唯一，与插入顺序无关。\quad 知识点：STL 容器（set/map）；类别：C++ 程序设计；难度：5；组别：提高级。}

\ans{8.}{C}
\expn{拓扑排序的定义要求：对每条有向边 \texttt{u -> v}，\texttt{u} 都排在 \texttt{v} 之前。拓扑序通常不唯一；含环的有向图不存在拓扑排序；拓扑排序与深度优先遍历序列不是一回事。\quad 知识点：拓扑排序；类别：图论；难度：6；组别：提高级。}

\ans{9.}{A}
\expn{\texttt{next} 数组（部分匹配表）记录模式串每个位置匹配失败后应回溯到的位置（基于最长公共前后缀），匹配过程中主串指针不回退。\quad 知识点：字符串匹配（KMP）；类别：算法；难度：6；组别：提高级。}

\ans{10.}{A}
\expn{朴素 Dijkstra 每轮用 \texttt{O(n)} 选择未确定的最小距离顶点、\texttt{O(n)} 松弛，共 \texttt{n} 轮，总复杂度 \texttt{O(n\^{}2)}。\quad 知识点：单源最短路（Dijkstra）；类别：图论；难度：6；组别：提高级。}

\ans{11.}{D}
\expn{模 20260913 下，\texttt{1/k + 1/(p-k)} 成对抵消（\texttt{1/(p-k) $\equiv$ -1/k}），故 \texttt{1/1+1/2+...+1/(p-1) $\equiv$ 0}；缺 \texttt{1/1=1}，和 $\equiv$ -1；指数 20260913 为奇数，\texttt{(-1)\^{}20260913 = -1 $\equiv$ 20260912}。\quad 知识点：模质数运算（倒数配对）；类别：数学与其他；难度：7；组别：提高级。}

\ans{12.}{C}
\expn{在 \texttt{[0,2]} 上独立均匀取两点，\texttt{E[|X-Y|]} 为单位区间情形的 2 倍，即 \texttt{2/3}。\quad 知识点：概率期望；类别：数学与其他}

\ans{13.}{D}
\expn{图中含节点数大于 1 的强连通分量恰有 2 个（两个三元环），其余顶点为平凡强连通分量。\quad 知识点：强连通分量；类别：图论；难度：7；组别：提高级。}

\ans{14.}{B}
\expn{容斥原理：\texttt{200-(100+66+28)+(33+14+9)-4 = 58}（\texttt{|2|=100}、\texttt{|3|=66}、\texttt{|7|=28}，两两交集 \texttt{33,14,9}，三交集 \texttt{4}）。\quad 知识点：容斥原理；类别：数学与其他；难度：7；组别：提高级。}

\ans{15.}{A}
\expn{满 \texttt{k} 叉树深度 \texttt{h} 的结点数为 \texttt{1+k+...+k\^{}(h-1) = (k\^{}h-1)/(k-1)}。\quad 知识点：树的性质（满 k 叉树）；类别：数据结构；难度：3；组别：入门级。}

\ans{16.}{A（√）}
\expn{函数 C 中 $k$ 为小于 $x$ 的元素个数，$k\cdot x-L$ 为左侧贡献，$P-L-(n-k)x$ 为右侧贡献，合起来恰为 $\sum|A[i]-x|$。\quad 知识点：曼哈顿距离与前缀和；类别：算法；难度：3；组别：入门级。}

\ans{17.}{B（×）}
\expn{扩大枚举范围后，新增的 $x$ 均满足 $\sum|A[i]-x|>D$，会被 \texttt{if (v <= D)} 过滤，不会进入 F 数组，故输出不变。\quad 知识点：枚举范围与可行性；类别：算法；难度：4；组别：入门级。}

\ans{18.}{A（√）}
\expn{当 $x<A[0]-D$ 时，$\sum|A[i]-x|\ge |A[0]-x|>D$，任何这样的 $x$ 都不满足条件，因此枚举从 $A[n-1]-D$ 到 $A[0]+D$ 已覆盖所有可行 $x$。\quad 知识点：曼哈顿距离下界估计；类别：算法；难度：4；组别：入门级。}

\ans{19.}{A}
\expn{\texttt{lower\_bound} 返回第一个不小于 $x$ 的迭代器，故 $k$ 为数组中小于 $x$ 的元素个数。\quad 知识点：二分查找（lower\_bound）；类别：算法；难度：4；组别：入门级。}

\ans{20.}{C}
\expn{$n=2$，$D=3$，点 $(0,0),(1,0)$。令 $f(x)=|x|+|x-1|$，需 $f(x)+2|y|\le 3$：$x=-1$ 时 $f=3$，仅 $y=0$；$x=0$ 或 $1$ 时 $f=1$，$y\in\{-1,0,1\}$；$x=2$ 时 $f=3$，仅 $y=0$。可行整数对共 $1+3+3+1=8$ 个。\quad 知识点：曼哈顿距离计数；类别：算法；难度：2；组别：入门级。}

\subsection*{（2）全局最小割}

\ans{21.}{A（√）}
\expn{该程序反复执行“最大权合并”（合并-缩点贪心策略），最终求得无向图全局最小割的边权和。\quad 知识点：全局最小割；类别：图论；难度：8；组别：提高级。}

\ans{22.}{A（√）}
\expn{$n=1$ 时 $ans$ 初始化为 0，主循环不执行，输出 0。\quad 知识点：边界情况；类别：算法；难度：8；组别：提高级。}

\ans{23.}{B（×）}
\expn{每轮循环结束时 \texttt{cur} 减 1，从 $n$ 降到 2 共执行 $n-1$ 轮（当 $cur=1$ 时循环条件不成立），并非 $n$ 轮。\quad 知识点：算法流程；类别：算法；难度：8；组别：提高级。}

\ans{24.}{B}
\expn{每轮按 $w$ 从大到小依次加入顶点，最后加入的顶点 $u$ 的 $w[u]$ 等于 $u$ 与已加入集合之间所有边的权值和，即当前划分下包含 $u$ 一侧的割大小，用于更新答案。\quad 知识点：最大权合并与割；类别：图论；难度：8；组别：提高级。}

\ans{25.}{B}
\expn{每轮内“选最大 $w$ 的顶点”与“更新 $w$”均为 $O(n\^{}2)$，共 $n-1$ 轮，总时间复杂度 $O(n\^{}3)$。\quad 知识点：时间复杂度分析；类别：算法；难度：8；组别：提高级。}

\ans{26.}{B}
\expn{运行程序或手工模拟得全局最小割为 4（割集为边 $(1,3)$、$(2,3)$、$(3,4)$，即顶点 3 与其余顶点之间的边，割边权和为 4）。\quad 知识点：全局最小割；类别：图论；难度：8；组别：提高级。}

\subsection*{（3）记忆值与离线分治}

\ans{27.}{A（√）}
\expn{程序将所有修改与查询统一记录为事件，再用分治离线计算每个查询的答案。\quad 知识点：分治（离线处理）；类别：算法。}

\ans{28.}{B（×）}
\expn{\texttt{add} 中 $x$ 沿 $x-=x\&-x$ 向左更新，\texttt{ask} 中 $x$ 沿 $x+=x\&-x$ 向右累加，因此 \texttt{ask(x)} 统计的是下标\textbf{不小于} $x$ 的树状数组元素之和，而非“不大于”。\quad 知识点：树状数组；类别：数据结构。}

\ans{29.}{A（√）}
\expn{某种形状只出现一次时其首次与最后出现位置相同，贡献为 0。\quad 知识点：记忆值定义；类别：算法。}

\ans{30.}{A}
\expn{\texttt{add\_e(u, v, w)} 中 \texttt{if (u \&\& v)} 表示仅当 $u$、$v$ 均非 0（相邻两次出现位置都存在）时才生成事件；$u=0$ 表示“哨兵”，不生成事件。\quad 知识点：事件构造；类别：算法。}

\ans{31.}{C}
\expn{分治共 $\log$ 层，每层内对事件排序归并，并在树状数组上做 $O(\log)$ 的修改/查询，总复杂度 $O((n+m)\log^2(n+m))$。\quad 知识点：分治复杂度；类别：算法。}

\ans{32.}{B}
\expn{初始序列 $1,2,3,1,3,2,1$，第一个查询为区间 $[3,7]$（元素 $3,1,3,2,1$）：形状 3 贡献 $5-3=2$，形状 1 贡献 $7-4=3$，形状 2 贡献 0，合计 $5$。\quad 知识点：记忆值计算；类别：算法。}

\clearpage
\section*{三、完善程序（33--42）}

\subsection*{（1）矩阵变换（P11717，稳定婚姻）}

\ans{33.}{A}
\expn{行 $i$ 按数字 $x$ 在行中位置 $p[i][x]$ 升序排序（越靠左越优先），故比较函数为 $p[i][x]<p[i][y]$。\quad 知识点：偏好序构造；类别：算法；难度：8；组别：提高级。}

\ans{34.}{B}
\expn{数字 $v$ 按 $p[i][v]$ 降序排序得到它对各行的偏好序，比较函数为 $p[i][v]>p[j][v]$。\quad 知识点：偏好序构造；类别：算法；难度：8；组别：提高级。}

\ans{35.}{A}
\expn{$rn[v][pn[v][k]]=k$ 记录数字 $v$ 对第 $pn[v][k]$ 行的偏好排名，供替换时 $O(1)$ 比较。\quad 知识点：排名数组；类别：算法；难度：8；组别：提高级。}

\ans{36.}{A}
\expn{数字 $v$ 已被行 $j$ 匹配，新来的行 $i$ 若在 $v$ 的偏好中排名更靠前（$rn[v][i]<rn[v][j]$），则 $v$ 改配 $i$、$j$ 被淘汰。\quad 知识点：稳定婚姻替换规则；类别：算法；难度：8；组别：提高级。}

\ans{37.}{B}
\expn{被淘汰的行 $j$ 重新入队（\texttt{q[t++] = j}），等待继续尝试下一个偏好数字，体现“不匹配者再次追求”。\quad 知识点：队列与贪心匹配；类别：算法；难度：8；组别：提高级。}

\subsection*{（2）树上咒语相撞（CF855G，树上计数 + 并查集）}

\ans{38.}{B}
\expn{DFS 回溯时累计子树大小 $sz[x] += sz[to[i]]$。\quad 知识点：树的 DFS 与子树大小；类别：算法；难度：6；组别：提高级。}

\ans{39.}{B}
\expn{$s[i]$ 为 $i$ 的各相邻分支大小的平方和，累加时使用 \texttt{1ll} 防止 int 溢出：$t += 1ll\cdot sz[to[j]]\cdot sz[to[j]]$。\quad 知识点：树形计数与防溢出；类别：算法；难度：7；组别：提高级。}

\ans{40.}{B}
\expn{并查集每个分量的“锚点”初始为单点自身，其对应原树中的子树大小为 $u[i]=sz[i]$。\quad 知识点：并查集初始化；类别：算法；难度：6；组别：提高级。}

\ans{41.}{A}
\expn{加边后需把深度大的端点沿父链逐层向上合并，故当 $dep[x]<dep[y]$ 时交换，保证 $x$ 为深度较大者。\quad 知识点：树深度与并查集合并；类别：算法；难度：6；组别：提高级。}

\ans{42.}{A}
\expn{合并两个分量时维护分量大小：$z[y] += z[x]$，同时 $f[x]=y$ 完成并查集合并。\quad 知识点：并查集按大小合并；类别：算法；难度：6；组别：提高级。}

\clearpage
\section*{附录：42 题考点统计表}
\begingroup
\small
\begin{longtable}{cp{2.2cm}p{2.0cm}ccp{5.0cm}}
\toprule
题号 & 知识点 & 类别 & 难度 & 组别 & 说明 \\
\midrule
1  & Linux 终端常用命令（cat） & 基础知识与编程环境 & 5 & 提高级 & cat 连接显示 \ \\
2  & 哈夫曼树性质 & 数据结构 & 4 & 入门级 & m 叶空指针 2m \ \\
3  & 完全二叉树与复杂度 & 算法 & 6 & 提高级 & 叶子路径 O(n log n) \ \\
4  & 卡特兰数 & 数学与其他 & 7 & 提高级 & (0,0)->(5,5) x>=y 42 \ \\
5  & 栈与队列的模拟 & 数据结构 & 3 & 入门级 & 双栈队列 \ \\
6  & 递推（斐波那契） & 算法 & 3 & 入门级 & 爬楼梯 f(10)=89 \\
7  & STL 容器（set/map） & C++ 程序设计 & 5 & 提高级 & 有序去重 O(log n) \\
8  & 拓扑排序 & 图论 & 6 & 提高级 & 边 u->v 则 u 在前 \\
9  & 字符串匹配（KMP） & 算法 & 6 & 提高级 & next 数组回溯 \\
10  & 单源最短路（Dijkstra） & 图论 & 6 & 提高级 & 朴素 O(n\^{}2) \ \\
11  & 模质数运算（倒数配对） & 数学与其他 & 7 & 提高级 & 20260912 \\
12  & 概率期望 & 数学与其他 & -- & -- & E|X-Y|=2/3 \ \\
13  & 强连通分量 & 图论 & 7 & 提高级 & 计数 2 \ \\
14  & 容斥原理 & 数学与其他 & 7 & 提高级 & 1~200 计数 58 \ \\
15  & 树的性质（满 k 叉树） & 数据结构 & 3 & 入门级 & (k\^{}h-1)/(k-1) \ \\
16  & 曼哈顿距离与前缀和 & 算法 & 3 & 入门级 & 阅读(1) 函数 C \ \\
17  & 枚举范围与可行性 & 算法 & 4 & 入门级 & 阅读(1) \ \\
18  & 曼哈顿距离下界估计 & 算法 & 4 & 入门级 & 阅读(1) \ \\
19  & 二分查找（lower\_bound） & 算法 & 4 & 入门级 & 阅读(1) \ \\
20  & 曼哈顿距离计数 & 算法 & 2 & 入门级 & 阅读(1) 样例 \ \\
21  & 全局最小割（合并-缩点） & 图论 & 8 & 提高级 & 阅读(2) \ \\
22  & 边界情况 & 算法 & 8 & 提高级 & 阅读(2) n=1 \ \\
23  & 算法流程 & 算法 & 8 & 提高级 & 阅读(2) 轮数 \ \\
24  & 最大权合并与割 & 图论 & 8 & 提高级 & 阅读(2) w[u] \ \\
25  & 时间复杂度分析 & 算法 & 8 & 提高级 & 阅读(2) O(n³) \ \\
26  & 全局最小割 & 图论 & 8 & 提高级 & 阅读(2) 样例 \ \\
27  & 分治（离线处理） & 算法 & -- & -- & 阅读(3) \ \\
28  & 树状数组 & 数据结构 & -- & -- & 阅读(3) ask 方向 \ \\
29  & 记忆值定义 & 算法 & -- & -- & 阅读(3) \ \\
30  & 事件构造 & 算法 & -- & -- & 阅读(3) add\_e \ \\
31  & 分治复杂度 & 算法 & -- & -- & 阅读(3) \ \\
32  & 记忆值计算 & 算法 & -- & -- & 阅读(3) 样例 \ \\
33  & 偏好序构造（行） & 算法 & 8 & 提高级 & 完善(1) \ \\
34  & 偏好序构造（列） & 算法 & 8 & 提高级 & 完善(1) \ \\
35  & 排名数组 & 算法 & 8 & 提高级 & 完善(1) \ \\
36  & 稳定婚姻替换规则 & 算法 & 8 & 提高级 & 完善(1) \ \\
37  & 队列与贪心匹配 & 算法 & 8 & 提高级 & 完善(1) \ \\
38  & 树的 DFS 与子树大小 & 算法 & 6 & 提高级 & 完善(2) \ \\
39  & 树形计数与防溢出 & 算法 & 7 & 提高级 & 完善(2) \ \\
40  & 并查集初始化 & 算法 & 6 & 提高级 & 完善(2) \ \\
41  & 树深度与并查集合并 & 算法 & 6 & 提高级 & 完善(2) \ \\
42  & 并查集按大小合并 & 算法 & 6 & 提高级 & 完善(2) \ \\
\bottomrule
\end{longtable}
\endgroup

\noindent\textbf{汇总：}
\begin{itemize}[leftmargin=2em, itemsep=1pt]
    \item 类别分布（42 题）：基础知识与编程环境 1（1）、C++ 程序设计 1（7）、数据结构 4（2/5/15/28）、算法 26、图论 6（8/10/13/21/24/26）、数学与其他 4（4/11/12/14）。\quad 说明：阅读程序与完善程序主体归入“算法”，图论算法题归入“图论”。
    \item 组别分布：入门级 9 题（2/5/6/15/16/17/18/19/20），提高级 26 题，另有 7 题（12、27--32）难度与组别未标注。
    \item 难度分布：2（1 题）、3（4）、4（4）、5（2）、6（8）、7（5）、8（11），共 42 题（其中 7 题难度未标注）。
    \item 考点去重：同一大题内知识点不重复；“输入给定数据求输出”型小题仅题 20、26、32 三处。
\end{itemize}

\end{document}
